\documentclass[10pt]{article}
\usepackage[margin=0.75in, top=0.65in, bottom=0.65in]{geometry}
\usepackage[compact]{titlesec}
\usepackage{listings}
\usepackage{xcolor}
\usepackage{enumitem}
\usepackage[hidelinks]{hyperref}

\setlength{\parindent}{0pt}
\setlength{\parskip}{4pt}
\titleformat{\section}[block]{\large\bfseries}{}{0pt}{\setlength{\leftskip}{0pt}}
\titleformat{\subsection}[block]{\normalsize\bfseries}{}{0pt}{\setlength{\leftskip}{0pt}}
\titlespacing{\section}{0pt}{8pt}{4pt}
\titlespacing{\subsection}{0pt}{6pt}{3pt}
\AtBeginDocument{\setlength{\leftskip}{1.5em}}

\lstset{
  basicstyle=\ttfamily\scriptsize,
  backgroundcolor=\color{gray!10},
  frame=single,
  breaklines=true,
  aboveskip=4pt,
  belowskip=4pt
}

\title{CMPT 276 Assignment 4}
\author{Krish Kaushik}
\date{March 31, 2026}

\begin{document}
\maketitle


% ─────────────────────────────────────────────────────────────────────────────
\section*{1. Code Smells Identified}
% ─────────────────────────────────────────────────────────────────────────────

I reviewed code implemented by \textbf{Tanveer Muntaseer} in \texttt{ObamaScreen.java} and \texttt{GameLogicHelper.java}, and identified the 4 code smells below.


\subsection*{Smell 1 - Dead Code \quad \normalfont\texttt{ObamaScreen.java}, \texttt{getProgressiveDialogue()}, lines 378--382}

\texttt{totalDuration} is assigned \texttt{TALK\_FALLBACK\_DURATION} unconditionally on line 378, then the \texttt{if (SettingsScreen.soundOn \&\& obamaVoice != null)} block immediately below assigns the exact same value again. The condition can never change the result, making the branch dead code that misleads future readers.

\textbf{Suggestion:} Delete the \texttt{if} block entirely.

\subsection*{Smell 2 - Inappropriate Intimacy \quad \normalfont\texttt{ObamaScreen.java}, lines 240, 251, 381}

\texttt{ObamaScreen} directly accesses \texttt{SettingsScreen.soundOn} in three separate places. These two classes have no architectural relationship (one is a cutscene, the other is a settings UI) creating tight coupling that makes refactoring either class riskier.

\textbf{Suggestion:} Add a \texttt{private boolean soundOn} field to \texttt{ObamaScreen}, pass it as a constructor parameter, and replace all three direct references with the local field so the coupling lives in exactly one place.

\subsection*{Smell 3 - Long Method \quad \normalfont\texttt{ObamaScreen.java}, \texttt{render()}, lines 230--266}

The \texttt{TALK} phase logic is handled entirely inline inside \texttt{render()}, while the other two phases already delegate to dedicated methods (\texttt{updateObamaWalk()}, \texttt{updatePlayerWalk()}). The inconsistency makes the method long and harder to read.

\textbf{Suggestion:} Extract the \texttt{TALK} case into a new private \texttt{updateTalkPhase(float delta)} method, so all three phases are handled the same way and \texttt{render()} becomes a clean coordinator.

\subsection*{Smell 4 - Data Clumps \quad \normalfont\texttt{GameLogicHelper.java}, lines 172--189}

\begin{lstlisting}
public static float getPlayerSpawnX(int mapIndex) { ... }
public static float getPlayerSpawnY(int mapIndex, float[] mapYOffsets) { ... }
\end{lstlisting}

Two parallel switch methods represent one concept: a spawn position. Every call site must invoke both and pass \texttt{mapIndex} twice. Adding a new map requires updating both switches in sync; forgetting one produces a wrong position with no compiler warning.

\textbf{Suggestion:} Merge into a single \texttt{getPlayerSpawnPosition(int mapIndex, float[] mapYOffsets)} returning a \texttt{Vector2}, so each map's coordinates are defined in one place.

% ─────────────────────────────────────────────────────────────────────────────
\section*{2: Refactoring My Own Code}
% ─────────────────────────────────────────────────────────────────────────────

After receiving a review from Samantha Gan, I addressed four code smells she identified in \texttt{PoliceEnemy.java} and \texttt{GameScreen.java}. Each fix was committed and pushed separately; \texttt{mvn test} was rerun after each commit and all tests passed every time.

\subsection*{Fix 1 - Duplicated Code \quad \normalfont\texttt{PoliceEnemy.java} \hfill commit \texttt{61d29a3}}

The barrier-collision logic -- computing \texttt{hitW}, \texttt{hitTop}, and four \texttt{isBarrier()} checks for both axes -- was copy-pasted verbatim in both the chasing branch and the wander branch of \texttt{update()}. I extracted it into a private \texttt{applyBarrierCollision(float newX, float newY, BarrierLookup lookup)} method (line 189) and replaced both duplicated blocks with a single call.

\subsection*{Fix 2 - Long Method: \texttt{update()} \quad \normalfont\texttt{PoliceEnemy.java} \hfill commit \texttt{2ccdb8a}}

\texttt{update()} handled the alert state machine, movement, and animation all inline across roughly 125 lines. I extracted three private helpers -- \texttt{updateAlertState}, \texttt{updateMovement}, and \texttt{updateAnimation} -- each with a single responsibility, reducing \texttt{update()} to four lines:

\begin{lstlisting}
float dist = Vector2.dst(position.x, position.y, playerX, playerY);
updateAlertState(delta, dist);
updateMovement(delta, playerX, playerY, lookup);
updateAnimation(delta);
\end{lstlisting}

\subsection*{Fix 3 - Dead Code \quad \normalfont\texttt{GameScreen.java} \hfill commit \texttt{cae4571}}

\texttt{activateMap()} called \texttt{police.clear()} then immediately checked \texttt{police.size() < 3} and \texttt{police.size() < 16} which is always true after a clear, making both guards dead code. When I implemented the clearing, I never changed the police size checks. To address this, I removed them and I also extracted the 33 hardcoded spawn calls into named constant 2D arrays (\texttt{MAP2\_POLICE\_SPAWNS}, \texttt{MAP3\_POLICE\_SPAWNS}) and a private \texttt{spawnEnemiesForMap(int index)} helper, so \texttt{activateMap()} delegates to a single call and spawn coordinates are defined once at the top of the class.

\subsection*{Fix 4 - Low Cohesion: \texttt{render()} \quad \normalfont\texttt{PoliceEnemy.java} \hfill commit \texttt{8711311}}

\texttt{render()} was doing two unrelated things: drawing the police sprite and drawing the alert indicator (\texttt{?} or \texttt{!}) overlay. I split it into two private methods: \texttt{renderSprite(batch)} and \texttt{renderAlertIndicator(batch)}. Now each has exactly one reason to change, and \texttt{render()} simply calls both.

\end{document}
